\chapter{Keys and mouse}\label{app:keys}
\begin{goals}
\textbf{This appendix} lists every keyboard shortcut and mouse gesture in wuffle, read from the program
itself. On a Mac, \key{⌘} works wherever \key{Ctrl} is given.
\end{goals}

\section{The desk}

\subsection{Keys}

\begin{center}\small
\begin{tabularx}{\linewidth}{@{}>{\raggedright}p{0.3\linewidth}L@{}}\toprule
\sffamily\bfseries key & \sffamily\bfseries what it does \\\midrule
\key{Ctrl}+\key{O} & open data, a project or a wuffle page (\cref{ch:data}) \\
\key{Ctrl}+\key{S} & save the project as \file{.zip.wuffle} (\cref{sec:project-file}) \\
\key{Ctrl}+\key{Z} & undo the last change, whoever made it (\cref{sec:undo}) \\
\key{Ctrl}+\key{Shift}+\key{Z}, \key{Ctrl}+\key{Y} & redo \\
\key{`} (backquote) & open the console (\cref{ch:console}) \\
\key{Esc} & clear the outlines of the selection shapes from the nets; again (or when there are none), the selection \\
\key{P} \key{L} \key{C} \key{B} \key{M} & the net's tools: pick, lasso, cone, band, measure (\cref{sec:select-tools}) \\
\key{R} & reset the active net's rotation, as its \ui{↺ rotated} button does (\cref{ch:rotation}) \\
\key{←} \key{→} \key{↑} \key{↓} & turn the active net: spin about its centre, tilt it; 5° a press, 1° with \key{Shift} \\
\key{Ctrl}+\key{A} & select everything the active net shows \\
\key{Ctrl}+\key{C} & copy the selected attitudes, one per line \\
\key{Ctrl}+\key{PgUp} / \key{PgDn} & the previous / next diagram \\
\key{?} & the keys and the mouse, in short \\
\bottomrule\end{tabularx}
\end{center}

The tool keys, \key{R}, the arrows, \key{Ctrl}+\key{A}, \key{Ctrl}+\key{C} and \key{?} act on the net or the bare page only; \menu{Tools} lists them too.
\menu{Help>Keyboard and mouse} shows the short version of this appendix. Undo, redo, \key{`} and \key{Esc} don't act while you type in a field, a dialog or the console, where
the keys belong to what you are typing. \key{Ctrl}+\key{O} and \key{Ctrl}+\key{S} work everywhere,
even there. Undo and redo also leave a figure tab alone: a figure has its own undo (below).

\subsection{On a net}

\begin{center}\small
\begin{tabularx}{\linewidth}{@{}>{\raggedright}p{0.3\linewidth}L@{}}\toprule
\sffamily\bfseries gesture & \sffamily\bfseries what it does \\\midrule
\multicolumn{2}{@{}l}{\sffamily\itshape with the \ui{pick} tool} \\
click a measurement & select its row, alone (\cref{sec:pick}) \\
\key{Shift}+click & add its row to the selection \\
\key{Ctrl}+click & add or remove its row \\
click empty space & clear the selection \\
drag & rotate the net (\cref{ch:rotation}); a click that moves under 4 pixels is still a click \\
\multicolumn{2}{@{}l}{\sffamily\itshape with \ui{lasso}, \ui{cone} or \ui{band}} \\
drag & draw the shape; the rows inside it become the selection (\cref{sec:shapes}) \\
\key{Shift}+drag & add the rows inside it \\
\key{Alt}+drag & take the rows inside it away \\
\multicolumn{2}{@{}l}{\sffamily\itshape always} \\
move the pointer & the status line reads the attitude under it, and the row under it if any \\
right-click & the net's menu: what is under the pointer, the selection, the net, the picture (\cref{sec:net-menu}); with \ui{measure}, a right \emph{drag} measures and a right \emph{click} opens the measure menu \\
double-click a measurement & its row in the table \\
hover, one measurement selected & the status line adds the angle from it \\
\bottomrule\end{tabularx}
\end{center}

\subsection{In the layers tree}

\begin{center}\small
\begin{tabularx}{\linewidth}{@{}>{\raggedright}p{0.3\linewidth}L@{}}\toprule
\sffamily\bfseries gesture & \sffamily\bfseries what it does \\\midrule
click the small arrow & open or close a diagram, a layer, a group \\
click a checkbox & show or hide a layer or one of its elements (\cref{ch:elements}) \\
click a layer or a dataset & show its rows in the table \\
double-click a layer & its properties (\cref{sec:properties}) \\
right-click & the menu of a diagram, a layer, a dataset or a figure \\
drag a layer onto a diagram's name & add a view of it there (\cref{sec:views}) \\
\key{Shift} while dropping a layer & move it there instead \\
drag a dataset onto a diagram's name & draw it there, as a new layer \\
\bottomrule\end{tabularx}
\end{center}

\subsection{Elsewhere on the desk}

\begin{description}
\item[Drop files on the window] to open them, as \menu{File>Open data or a project…} does: data, a
\file{.zip.wuffle}, a wuffle page, or a field log with its \file{.sig} and photos together.
\item[Drag a tab] by its title to put it beside, below or above the others; the desk's regions are
tabs you arrange.
\item[In the table,] click a row, or drag down several, to select them everywhere; the arrow keys,
\key{Home}, \key{End} and \key{Page Up}/\key{Page Down} move through it, and \key{Shift} extends a range.
\end{description}

\section{The console}

\begin{center}\small
\begin{tabularx}{\linewidth}{@{}>{\raggedright}p{0.3\linewidth}L@{}}\toprule
\sffamily\bfseries key & \sffamily\bfseries what it does \\\midrule
\key{Enter} & run the line \\
\key{Tab} & complete: a command at the start of the line, then a name, then a path \\
\key{↑} \key{↓} & earlier and later lines \\
\key{Backspace} & delete the last character \\
\key{Ctrl}+\key{C} & abandon the line; a fresh prompt \\
\bottomrule\end{tabularx}
\end{center}

The line editor works at the end of the line only: \key{←} and \key{→} don't move the cursor. Undo and
redo are the commands \code{undo} and \code{redo} here.

\section{A figure tab}

A figure is a page of \code{@gcu/compo} frames (\cref{ch:figures}), with an editor of its own:

\begin{center}\small
\begin{tabularx}{\linewidth}{@{}>{\raggedright}p{0.3\linewidth}L@{}}\toprule
\sffamily\bfseries key or gesture & \sffamily\bfseries what it does \\\midrule
click a frame & select it; \key{Shift}+click adds to the selection \\
drag on the empty page & select the frames the box touches \\
drag a frame & move it, snapping to the others' edges \\
drag a handle & resize \\
arrow keys & nudge the selected frames; with \key{Shift}, ten times as far. A run of nudges is one undo step \\
\key{Delete}, \key{Backspace} & delete the selected frames (locked ones stay) \\
\key{Ctrl}+\key{D} & duplicate \\
\key{Ctrl}+\key{C}, \key{Ctrl}+\key{X}, \key{Ctrl}+\key{V} & copy, cut and paste frames, also from one figure to another \\
\key{Ctrl}+\key{Alt}+\key{C}, \key{Ctrl}+\key{Alt}+\key{V} & copy one frame's style (border, fill, type), put it on the selected frames \\
\key{Ctrl}+\key{A} & select every frame \\
\key{Ctrl}+\key{G}, \key{Ctrl}+\key{Shift}+\key{G} & group, ungroup \\
\key{Ctrl}+\key{]}, \key{Ctrl}+\key{[} & one step forward, back; with \key{Shift}, to the front, to the back \\
\key{Ctrl}+\key{Z}, \key{Ctrl}+\key{Shift}+\key{Z} or \key{Ctrl}+\key{Y} & the figure's own undo and redo \\
\key{Esc} & select nothing \\
\bottomrule\end{tabularx}
\end{center}

\section{The compass}

\begin{center}\small
\begin{tabularx}{\linewidth}{@{}>{\raggedright}p{0.3\linewidth}L@{}}\toprule
\sffamily\bfseries key or gesture & \sffamily\bfseries what it does \\\midrule
\ui{plane}, \ui{line} & what the next measurement is (\cref{sec:measuring}) \\
\ui{◎} & keep the live reading \\
\key{Enter} in the typed field, or \ui{add} & add the typed measurement (\cref{sec:manual-entry}) \\
\ui{×} on a row & delete that measurement \\
\ui{photo}, then \ui{capture} or \ui{cancel} & photograph the outcrop for the last measurement (\cref{sec:photos}) \\
\ui{clear} & empty the log \\
\ui{publish →} & write the log out (\cref{sec:publish}) \\
\bottomrule\end{tabularx}
\end{center}

\begin{caution}
\ui{clear} asks first: one tap says \ui{tap clear again to delete all …}, and only a second tap within
three seconds empties the log. There is no undo after that. Publish first.
\end{caution}
